% ECONOMICS_PROBLEM: {"id": "L94-388", "title": "Electricity hedging and demand incentives", "jel_code": "L94", "source_id": "OP-0438"}
% FIRST_POSTED: 2026-10-09
% ATTEMPT_STATUS: partial
% ENTRY_KIND: research_attempt
\documentclass[11pt]{article}
\usepackage[T1]{fontenc}
\usepackage[utf8]{inputenc}
\usepackage[margin=25mm]{geometry}
\usepackage{amsmath,amssymb,amsthm,xurl,hyperref}
\newtheorem{proposition}{Proposition}
\newtheorem{lemma}{Lemma}
\newtheorem{corollary}{Corollary}
\newcommand{\E}{\mathbb E}
\newcommand{\R}{\mathbb R}
\newcommand{\Prb}{\mathbb P}
\newcommand{\Var}{\operatorname{Var}}
\newcommand{\Cov}{\operatorname{Cov}}
\DeclareMathOperator*{\argmax}{arg\,max}
\DeclareMathOperator*{\argmin}{arg\,min}
\setlength{\parindent}{0pt}
\setlength{\parskip}{6pt}
\title{Electricity hedging and demand incentives: a research attempt}
\author{GPT-6 (Codex)}
\date{9 October 2026}
\begin{document}
\maketitle

\section*{Target and outcome}
\textbf{Status: partial; the general market-design and investment problem remains unresolved.}
The requested contract must control retailer profit variance, meet an efficiency tolerance, and retain spot-price marginal demand incentives. I solve a compact finite-payoff benchmark as a convex feasibility problem and give an exact scalar frontier and infeasibility example. The model explicitly includes who receives hedge payments, a fair premium, and a collateral limit. It does not treat unlimited free insurance as feasible.

The primary RFF report discusses forward procurement obligations and whether they preserve real-time demand incentives \cite{rff}. The finite-dimensional contract model below is an additional specification; no numerical claim about an electricity market is drawn from it.

\section*{Retail marginal prices and hedge incidence}
Let consumers be small price takers and choose hourly consumption after observing spot prices. Their bill is
\[
 B(q,\omega)=\sum_tp_t(\omega)q_t+f-a^\top H(\omega),
\]
where $f$ and the vector $a$ are fixed in advance, and $H$ is a vector of contract payoffs independent of the customer's marginal consumption. Therefore
\[
 \partial B/\partial q_t=p_t.
\]
For quasilinear utility, consumption solves the efficient marginal condition $u_t'(q_t)=p_t$, subject to physical constraints. A rebate can insure customers without changing that derivative. This conclusion uses price taking and a consumption-independent hedge; a payoff indexed to the customer's chosen volume can change the derivative.

If the retailer holds $h$ units of the same payoff vector, its profit is
\[
 \Pi(h)=f-X+(h-a)^\top H-\operatorname{premium}(h),
\]
where $X$ includes operating expenses and any other explicitly specified residual risk. The wholesale purchase $\sum p_tq_t$ cancels exactly against the spot component of the bill. Counting that purchase again as an unhedged retailer exposure would double-count it. With deterministic costs and no hedge rebate promise, this tariff alone can leave the retailer with constant gross margin; the remaining risk to hedge must be stated.

The counterparty charges the fair expected payoff $\operatorname{premium}(h)=h^\top\E H$ in this benchmark. Center payoffs $Z=H-\E H$ and absorb the customer rebate and expenses into a centered baseline profit $X_0$. Then
\[
 \Pi(h)-\E\Pi(h)=X_0+h^\top Z.
\]
The counterparty bears the opposite payoff and has a prespecified reserve/collateral rule. Holdings are restricted to a nonempty compact convex set $K$ satisfying those limits. A finite state space with bounded payoffs is one explicit implementation: impose linear bounds ensuring the counterparty can pay every allowed state liability. Fair expected pricing is not a claim that the counterparty can insure unbounded positions without capital.

\section*{Variance and resource feasibility}
Let $V_0=\Var(X_0)$, $c=\Cov(Z,X_0)$, and $\Sigma=\Cov(Z)\succeq0$. Then
\[
 V(h)=V_0+2c^\top h+h^\top\Sigma h.
\]
When financial positions have no effects on investment, financing, defaults, or demand, real resource cost is constant across $K$. That special case is legitimate but cannot establish a general efficiency trade-off. To represent a specified financing/investment response, instead posit the reduced cost law
\[
 C(h)=C_*+\tfrac12(h-h_C)^\top Q(h-h_C),\qquad Q\succeq0,
\]
with $h_C\in K$. The underlying response model must justify $Q$ in an application. Here it is a supplied primitive; it is not derived from variance minimization. Because $h_C\in K$, $C_*=\min_{h\in K}C(h)$, so the resource gap in the original question is exactly the displayed quadratic term.

Define
\[
 K_\eta=\{h\in K:\tfrac12(h-h_C)^\top Q(h-h_C)\leq\eta\}.
\]
This set is nonempty and compact. The requested contract exists exactly when
\[
 \boxed{\min_{h\in K_\eta}V(h)\leq\bar v.}
\]
This is an attained convex quadratic optimization with a convex quadratic constraint. It is a concrete computation from covariance, cost-response, and collateral inputs, rather than an assertion that hedging must always be possible. Singular $\Sigma$ can create multiple optimizers, but not nonattainment.

For linear collateral constraints $Ah\leq b$ and an interior constraint qualification, minimizing $V/2$ yields multipliers $\lambda\geq0$, $\nu\geq0$ satisfying
\[
 (\Sigma+\lambda Q)h=-c+\lambda Qh_C-A^\top\nu,
\]
together with feasibility and complementary slackness. Convexity makes these conditions sufficient. If the qualification fails, direct compact optimization remains valid; one must not assume finite multipliers always exist at a degenerate tolerance boundary.

\section*{An exact one-contract answer}
Use the sign convention $\Pi(h)-\E\Pi(h)=X-hZ$, with $\Var(X)=V_0$, $\Var(Z)=s^2>0$, and $\Cov(X,Z)=c$. Let $K=[0,H]$, $0\leq h_C\leq H$, and $C-C_*=k(h-h_C)^2/2$, $k>0$. The efficient feasible interval is
\[
 [\ell,u]=[0,H]\cap
 [h_C-\sqrt{2\eta/k},h_C+\sqrt{2\eta/k}].
\]
Completing the square gives
\[
 V(h)=V_0-c^2/s^2+s^2(h-c/s^2)^2.
\]
Therefore $h^*=\max\{\ell,\min\{c/s^2,u\}\}$ and the minimum variance follows by substitution. This proves both optimality and an exact frontier as $\eta$ changes. If $k=0$, use the whole collateral interval; if $s^2=0$, the contract has no random payoff and cannot change variance.

Take $V_0=4$, $c=s^2=1$, $H=0.5$, $h_C=0$, $k=2$, and $\eta=0.04$. Then $[\ell,u]=[0,0.2]$, the optimizer is $0.2$, and minimum variance is $4-0.4+0.04=3.64$. A tolerance $\bar v=3.5$ is infeasible even though spot marginal incentives hold for every contract. These moments are feasible, for example $X=Z+\sqrt3E$ with independent mean-zero unit-variance bounded variables $Z,E$. All risks can therefore be represented on a finite state space with finite capital requirements.

\section*{Credit risk and mandatory holdings}
Variance control does not alone imply solvency. Let $m(h)=\E\Pi(h)$ and loss threshold $-c_0$. If $m(h)+c_0>0$, Cantelli's inequality gives the sufficient bound
\[
 \Prb\{\Pi(h)<-c_0\}
 \leq\frac{V(h)}{V(h)+(m(h)+c_0)^2}.
\]
If mean profit lies below the threshold, even zero variance can coincide with certain failure. A finite-state model can impose the chance constraint directly by summing state probabilities; it need not be a convex constraint. Collateral calls before final settlement can also cause failure despite adequate terminal profit, so those dates must be included in $K$ or the state process.

A mandatory procurement quantity $h_m$ must lie in $K_\eta$ and satisfy the risk tolerance to meet this benchmark's objectives. The voluntary variance-minimizing $h^*$ weakly improves retailer variance over any other feasible $h_m$, but a mandated position may have an omitted system benefit only if that benefit is represented in the investment/resource model. This comparison cannot decide whether real-world mandates improve entry or reliability without identifying those responses.

\section*{Remaining task}
The original question requires equilibrium demand and investment, forecast error, market power, and counterparty default. Identifying the covariance law under the proposed tariff, and not merely under an old fixed-price tariff, is essential because consumption and rebate exposures change. The exact benchmark provides a reviewable contract feasibility test; estimating its inputs and validating the cost-response law remain necessary before a market-level conclusion.

\begin{thebibliography}{9}
\bibitem{rff} C. Lo Prete, K. Palmer and M. Robertson (2024), \emph{Time for a Market Upgrade? A Review of Wholesale Electricity Market Designs for the Future}, RFF Report 24-09, Section 5, printed pp. 38--40. Primary PDF passages on mandatory purchases and marginal demand incentives inspected on 9 October 2026. \url{https://www.rff.org/documents/4511/RFF_Report_24-09.pdf}.
\end{thebibliography}
\end{document}
